\begin{table}[htbp]
\centering
\footnotesize
\caption{\textbf{6.3} AUC de cada dimensi\'on, con los valores crudos (R con el signo invertido), y de la suma de los tres puntajes, como clasificadores de compra$_i$.}
\label{tab:p6_auc}
\begin{tabular}{lc}
\toprule
\textbf{clasificador} & \textbf{AUC} \\
\midrule
R (signo invertido) & 0.6764 \\
F (crudo) & 0.7210 \\
M (crudo) & 0.7093 \\
suma de los tres puntajes & 0.7352 \\
\bottomrule
\end{tabular}
\\[0.3em]{\scriptsize \textbf{Qu\'e leer:} la suma le gana a las tres por separado, que es la justificaci\'on emp\'irica de usar tres n\'umeros y no uno. La mejora es modesta porque las tres dimensiones est\'an correlacionadas entre s\'i.}
\end{table}
